\documentclass[11pt,a4paper]{article}

\usepackage[T1]{fontenc}
\usepackage[utf8]{inputenc}
\usepackage[spanish]{babel}
\usepackage{geometry}
\usepackage{microtype}
\usepackage{setspace}
\usepackage{booktabs}
\usepackage{tabularx}
\usepackage{enumitem}
\usepackage{hyperref}
\usepackage{xcolor}
\usepackage{titlesec}
\usepackage{fancyhdr}
\usepackage{array}
\usepackage{parskip}
\usepackage{tikz}

\usetikzlibrary{arrows.meta,positioning,fit,backgrounds}

\geometry{margin=2.3cm, top=2.6cm, bottom=2.6cm}
\setstretch{1.15}
\setlist[itemize]{leftmargin=1.4em,itemsep=0.28em,topsep=0.4em}
\setlist[enumerate]{leftmargin=1.4em,itemsep=0.28em,topsep=0.4em}

\definecolor{upnavy}{HTML}{1F3A5F}
\definecolor{upgray}{HTML}{5A6570}
\definecolor{upsoft}{HTML}{E8EEF4}

\hypersetup{colorlinks=true,urlcolor=upnavy,linkcolor=upnavy}
\titleformat{\section}{\large\bfseries\color{upnavy}}{\thesection.}{0.6em}{}
\titleformat{\subsection}{\normalsize\bfseries\color{upnavy}}{\thesubsection.}{0.5em}{}
\titleformat{\subsubsection}{\normalsize\itshape\color{upnavy}}{\thesubsubsection.}{0.5em}{}

\pagestyle{fancy}
\fancyhf{}
\fancyhead[L]{\small\color{upgray}Pipeline local --- defensoria\_extract}
\fancyhead[R]{\small\color{upgray}Universidad del Pacífico}
\fancyfoot[C]{\thepage}
\renewcommand{\headrulewidth}{0.3pt}

\newcolumntype{Y}{>{\raggedright\arraybackslash}X}

\begin{document}

\begin{titlepage}
\centering
\vspace*{1.4cm}
{\color{upnavy}\rule{\textwidth}{1.4pt}}\\[0.7em]
{\LARGE\bfseries\color{upnavy} El pipeline local:}\\[0.35em]
{\LARGE\bfseries\color{upnavy} de PDF a un panel caso $\times$ mes}\\[0.7em]
{\Large Extractor \texttt{defensoria\_extract}\\[0.2em]
(directorio Asistente)}\\[0.7em]
{\color{upnavy}\rule{\textwidth}{1.4pt}}\\[1.3em]
{\large Universo de conflictos de la Defensoría,\\[0.2em]
identidad por embeddings y capas de mejora}\\[2em]

\begin{flushleft}
\renewcommand{\arraystretch}{1.35}
\begin{tabular}{@{}l@{\hspace{1.2em}}l@{}}
\textbf{Proyecto} & Minería, conflictos sociales e indicadores territoriales \\
\textbf{Institución} & Universidad del Pacífico \\
\textbf{Fuente} & Reportes mensuales de conflictos sociales, Defensoría del Pueblo \\
\textbf{Código} & \texttt{02\_codigo/extraer\_conflictos.py} + \texttt{defensoria\_extract/} \\
\textbf{Salidas} & \texttt{panel\_caso\_mes.csv}; \texttt{casos\_unicos.csv} \\
\textbf{Cobertura} & 269 ediciones (n.$^\circ$~1--269, abril 2004 -- julio 2026) \\
\textbf{Identidad} & Embeddings (umbral 0{,}90); fuzzy opcional \\
\textbf{Fecha de esta nota} & Septiembre 2026 \\
\end{tabular}
\end{flushleft}

\vfill
\begin{flushleft}
\color{upgray}\small
Bitácora del pipeline \emph{nuevo}: cómo se lee el PDF nativo, cómo se arma el universo, cómo funciona la identidad por embeddings y qué se puede mejorar a mano o con Gemini. No es un paper y no se compiló a propósito; vive en \texttt{04\_docs}. El pipeline viejo de Drive está en \texttt{pipeline\_drive\_conflictos.tex}; el contraste largo, en \texttt{metodologia\_extraccion\_conflictos.tex}.
\end{flushleft}
\end{titlepage}

\tableofcontents
\newpage

%----------------------------------------------------------------------
\section{En una frase}

El producto del extractor local es el \textbf{universo}: todas las fichas que se pueden leer de los reportes mensuales, mes a mes, con minería como \emph{bandera} y no como filtro que borra el resto. Encima se pueden construir recortes (solo socioambiental, solo mención minera, catálogo de minas del proyecto), resúmenes con Gemini y tablas curadas. La extracción no decide por el analista qué conflictos ``cuentan''.

Corrida de referencia: 269 PDF, del orden de 35\,000 filas caso $\times$ mes, $\approx$4\,000 \texttt{caso\_id} con embeddings a umbral 0{,}90. Falta solo el PDF de junio 2005.

%----------------------------------------------------------------------
\section{Flujograma actual del proyecto}

Diagrama del flujo vigente (directorio Asistente). El núcleo es determinista; Gemini, el notebook y la QA son capas opcionales encima del panel. Las cajas rellenas son el flujo principal; las blancas, lo opcional.

\begin{center}
\begin{tikzpicture}[
  font=\footnotesize,
  node distance=0.38cm and 0.9cm,
  box/.style={draw=upnavy, rounded corners=2pt, align=center, inner sep=5pt,
              fill=upsoft, text=upnavy, text width=8.4cm, minimum height=0.58cm},
  sub/.style={draw=upnavy, rounded corners=2pt, align=left, inner sep=5pt,
              fill=white, text=upnavy, text width=8.4cm, minimum height=0.55cm},
  opt/.style={draw=upgray, rounded corners=2pt, align=center, inner sep=4pt,
              fill=white, text=upnavy, text width=3.5cm, minimum height=0.62cm},
  arr/.style={-{Latex}, thick, color=upnavy},
  arro/.style={-{Latex}, thick, color=upgray, dashed},
  stage/.style={draw=upgray!60, rounded corners=3pt, inner sep=6pt, fill=upsoft!35}
]
% ----- A. Corpus -----
\node[box] (a1) {A1.\ Descarga web (\texttt{descargar\_reportes\_defensoria.py})};
\node[box, below=of a1] (a2) {A2.\ Completar faltantes desde Drive (\texttt{completar\_faltantes\ldots})};
\node[box, below=of a2] (a3) {A3.\ Consolidar en \texttt{01\_datos/crudos/pdfs} (269 mensuales; falta jun.\ 2005)};

% ----- B. Inventario -----
\node[box, below=0.55cm of a3] (b1) {B1.\ Fechas: n.$^\circ$, mes, año desde portada (\texttt{extraer\_fechas\_reportes.py})};
\node[box, below=of b1] (b2) {B2.\ Formato de maqueta: 6 tipos + especial (\texttt{clasificar\_formatos\_pdf.py})};
\node[sub, below=of b2] (b3) {B3.\ Decide modo de lectura: texto corrido (2004--08) o dos columnas (2008--26)};

% ----- C. Extracción -----
\node[box, below=0.55cm of b3] (c1) {C1.\ CLI: \texttt{python -u 02\_codigo/extraer\_conflictos.py}};
\node[sub, below=of c1] (c2) {C2.\ \texttt{layout.py} --- PyMuPDF: bloques con $(x,y)$; columna L/R según $x$ medio};
\node[sub, below=of c2] (c3) {C3.\ \texttt{secciones.py} --- inicio/fin del detalle; estado Activo/Latente/\ldots};
\node[sub, below=of c3] (c4) {C4.\ \texttt{fichas.py} --- abre/cierra ficha (puede cruzar página); parsea Tipo/Caso/\ldots};
\node[sub, below=of c4] (c5) {C5.\ Limpieza de basura de maquetación; completa departamento si falta};

% ----- D. Banderas + identidad -----
\node[box, below=0.55cm of c5] (d1) {D1.\ \texttt{flags.py} --- \texttt{es\_socioambiental}, \texttt{menciona\_mineria} (no borra filas)};
\node[sub, below=of d1] (d2) {D2.\ Embeddings: estandarizar texto $\rightarrow$ vector MiniLM multilingual};
\node[sub, below=of d2] (d3) {D3.\ Clustering por departamento, umbral coseno 0{,}90 $\rightarrow$ \texttt{caso\_id}};
\node[sub, below=of d3] (d4) {D4.\ (Alt.) \texttt{--identidad fuzzy}: rapidfuzz umbral 88, también por depto.};

% ----- E. Producto -----
\node[box, below=0.55cm of d4] (e1) {E1.\ \texttt{panel\_caso\_mes.csv} --- una fila = un caso en un mes ($\sim$35\,000)};
\node[box, below=of e1] (e2) {E2.\ \texttt{casos\_unicos.csv} --- un \texttt{caso\_id} + \texttt{n\_meses} ($\sim$4\,000)};
\node[box, below=of e2] (e3) {E3.\ \texttt{fallos.csv} + \texttt{RESUMEN.txt}};

\draw[arr] (a1) -- (a2);
\draw[arr] (a2) -- (a3);
\draw[arr] (a3) -- (b1);
\draw[arr] (b1) -- (b2);
\draw[arr] (b2) -- (b3);
\draw[arr] (b3) -- (c1);
\draw[arr] (c1) -- (c2);
\draw[arr] (c2) -- (c3);
\draw[arr] (c3) -- (c4);
\draw[arr] (c4) -- (c5);
\draw[arr] (c5) -- (d1);
\draw[arr] (d1) -- (d2);
\draw[arr] (d2) -- (d3);
\draw[arr] (d3) -- (d4);
\draw[arr] (d4) -- (e1);
\draw[arr] (e1) -- (e2);
\draw[arr] (e2) -- (e3);

% Marcos de etapa (detrás)
\begin{scope}[on background layer]
  \node[stage, fit=(a1)(a2)(a3), label={[font=\scriptsize\bfseries\color{upnavy}]above left:A.\ Corpus}] {};
  \node[stage, fit=(b1)(b2)(b3), label={[font=\scriptsize\bfseries\color{upnavy}]above left:B.\ Inventario}] {};
  \node[stage, fit=(c1)(c2)(c3)(c4)(c5), label={[font=\scriptsize\bfseries\color{upnavy}]above left:C.\ Extracción}] {};
  \node[stage, fit=(d1)(d2)(d3)(d4), label={[font=\scriptsize\bfseries\color{upnavy}]above left:D.\ Banderas e identidad}] {};
  \node[stage, fit=(e1)(e2)(e3), label={[font=\scriptsize\bfseries\color{upnavy}]above left:E.\ Producto}] {};
\end{scope}

% ----- F. Opcional (derecha de E) -----
\node[opt, right=1.15cm of e1] (f1) {F1.\ Notebook\\leer panel};
\node[opt, below=0.35cm of f1] (f2) {F2.\ Gemini\\(Vertex o API key)};
\node[opt, below=0.35cm of f2] (f3) {F3.\ Muestra dorada\\+ \texttt{evaluar\_extraccion}};
\node[opt, below=0.35cm of f3] (f4) {F4.\ \texttt{comparar\_identidad}};
\node[opt, below=0.35cm of f4] (f5) {F5.\ Cruce Drive 83\\\texttt{recrear\_tabla\ldots}};

\draw[arro] (e1.east) -- (f1.west);
\draw[arro] (e1.east) |- (f2.west);
\draw[arro] (e2.east) |- (f3.west);
\draw[arro] (e2.east) |- (f4.west);
\draw[arro] (e3.east) |- (f5.west);
\end{tikzpicture}
\end{center}

\begin{center}
\small\color{upgray}
Relleno = paso del núcleo.\quad
Borde blanco = detalle interno o capa opcional.\quad
Flecha punteada = no modifica el panel crudo.
\end{center}

\subsection{Lectura del diagrama}

\begin{enumerate}
  \item \textbf{A. Corpus.} Unir descarga web y PDF de Drive en una sola carpeta; excluir especiales y mal nombrados.
  \item \textbf{B. Inventario.} Fechas y etiqueta de maqueta; eso fija si se parsea en columnas o en texto corrido.
  \item \textbf{C. Extracción.} Un CLI; PyMuPDF con geometría; detector de sección; fichas con campos anclados; sin iLovePDF.
  \item \textbf{D. Banderas e identidad.} Minería como bit; embeddings (default) o fuzzy; \texttt{caso\_id} = cluster por departamento.
  \item \textbf{E. Producto.} Panel universo, casos únicos, fallos y resumen de corrida.
  \item \textbf{F. Opcional.} Exploración en notebook; resúmenes Gemini; QA manual; calibrar identidad; cruzar con los 83 de Drive.
\end{enumerate}

Comando central:

\begin{verbatim}
python -u 02_codigo/extraer_conflictos.py
\end{verbatim}

Opciones útiles: \texttt{--max N}, \texttt{--solo-numero 169}, \texttt{--identidad fuzzy}, \texttt{--umbral 0.90}. Rutas en \texttt{02\_codigo/rutas.py}. No hace falta iLovePDF.

%----------------------------------------------------------------------
\section{Cómo funciona el pipeline}

\subsection{Corpus e inventario}

Los PDF viven en \texttt{01\_datos/crudos/pdfs} (unión de descarga web + lo que había en Drive Minería). Scripts de apoyo: \texttt{descargar\_reportes\_defensoria.py}, \texttt{consolidar\_pdfs\_defensoria.py}, \texttt{completar\_faltantes\_defensoria.py}.

Antes de parsear fichas:

\begin{enumerate}
  \item \texttt{extraer\_fechas\_reportes.py} --- número de edición, mes y año (portada primero; nombre de archivo de respaldo).
  \item \texttt{clasificar\_formatos\_pdf.py} --- etiqueta de maqueta por PDF (lista narrativa 2004, tabla por regiones, unidad, sumilla, Adjuntía, infografía 2024--26). Eso decide si la lectura es \textbf{texto corrido} o \textbf{dos columnas}.
\end{enumerate}

Si hay dos PDF del mismo número, se queda el de más páginas. Se saltan especiales y el archivo mal nombrado.

\subsection{Lectura geométrica del PDF}

Paquete \texttt{defensoria\_extract}, motor \textbf{PyMuPDF} (\texttt{fitz}):

\begin{itemize}
  \item \texttt{layout.py} --- bloques con bbox $(x_0,y_0,x_1,y_1)$; columna izquierda si el punto medio $x<300$, derecha si no; limpia cabeceras y pies.
  \item \texttt{secciones.py} --- detector de inicio/fin del detalle de fichas (con guardas para no abrir en la sumilla); estado Activo / Latente / \ldots desde títulos de sección.
  \item \texttt{fichas.py} --- abre/cierra fichas (pueden cruzar páginas); parsea campos con regex anclada al inicio de línea (\texttt{Tipo:}, \texttt{Caso:}, ubicación, actores, etc.).
\end{itemize}

Diferencia de método respecto a Drive: no se entrega el layout a un conversor. Hechos del mes van a la derecha porque \emph{están} a la derecha.

Limitación explícita: desde $\approx$2018 los latentes en tablas compactas \textbf{sin} \texttt{Tipo:} no salen como filas.

\subsection{Banderas, no filtro}

\texttt{flags.py} marca sobre el texto de la ficha:

\begin{itemize}
  \item \texttt{es\_socioambiental} --- el tipo contiene ``socioambiental''.
  \item \texttt{menciona\_mineria} --- léxico (mina, relave, tajo, unidades conocidas, hidrocarburos del léxico, metales como palabra completa).
\end{itemize}

Nada se borra. Para análisis minero: filtrar \texttt{es\_socioambiental=1} y/o \texttt{menciona\_mineria=1}.

\subsection{Salidas}

\begin{center}
\renewcommand{\arraystretch}{1.15}
\begin{tabularx}{\textwidth}{@{}l Y@{}}
\toprule
\textbf{Archivo} & \textbf{Contenido} \\
\midrule
\texttt{panel\_caso\_mes.csv} & Una fila = un caso en un mes (estado, depto, tipo, caso, actores, hechos, banderas, formato, página, texto crudo). \\
\texttt{casos\_unicos.csv} & Una fila = un \texttt{caso\_id} (texto, \texttt{n\_meses}, banderas OR en el tiempo). \\
\texttt{fallos.csv} & PDF que no abrieron o salieron a cero fichas. \\
\texttt{RESUMEN.txt} & Conteos por formato y por reporte. \\
\bottomrule
\end{tabularx}
\end{center}

Herramientas del núcleo: PyMuPDF, pandas/csv, rapidfuzz (alternativa), \texttt{sentence-transformers}, scikit-learn, NumPy.

%----------------------------------------------------------------------
\section{Identidad por embeddings (en detalle)}

Esta es la decisión que más afecta cuántos ``conflictos'' ve el analista. La Defensoría \textbf{reescribe} el mismo caso mes a mes. Sin agrupar, cada redacción distinta sería un caso nuevo. Con agrupar de más, dos conflictos distintos del mismo departamento se vuelven uno.

\subsection{Qué problema resuelve}

El método fuzzy (\texttt{rapidfuzz}, \texttt{token\_sort\_ratio}) compara \emph{cómo está escrito} el texto. Si un año dice ``pobladores de Espinar exigen \ldots Antapaccay'' y otro alarga la ficha con la mesa de 2019, el string ya no se parece lo suficiente y el mismo conflicto se \textbf{fragmenta} en varios \texttt{caso\_id}.

Los \textbf{embeddings} proyectan cada descripción a un vector de significado. Dos redacciones del mismo conflicto quedan cerca aunque cambien las palabras. Eso tolera la reescritura de la Defensoría mejor que el fuzzy puro.

\subsection{Qué es un cluster / \texttt{caso\_id}}

Un \textbf{cluster} es el conjunto de fichas (filas del panel) que el algoritmo trata como el \textbf{mismo conflicto}. Ese montón recibe un solo \texttt{caso\_id}.

Ejemplo: Espinar--Antapaccay, ``cumplimiento de acuerdos de la mesa de 2013'', aparece en decenas de reportes con texto casi igual $\rightarrow$ un cluster (p.\ ej.\ \texttt{caso\_id}=1001) con \texttt{n\_meses} alto. Si en algunos meses la ficha enfatiza otras comunidades o el proyecto Coroccohuayco, puede caer en \textbf{otro} cluster aunque en la realidad sea el mismo conflicto Espinar. Por eso \texttt{n\_meses} cuenta apariciones de \emph{ese} cluster, no un reloj político perfecto.

\subsection{Pasos del algoritmo (\texttt{identidad\_emb.py})}

\begin{enumerate}
  \item \textbf{Estandarizar} el campo \texttt{caso}: sin tildes, minúsculas, sin stopwords (\texttt{estandarizar} en \texttt{identidad.py}). Queda \texttt{caso\_estandarizado}.
  \item \textbf{Deduplicar} dentro de cada departamento: la misma descripción repetida muchos meses se codifica \textbf{una sola vez} (se usa el texto más largo como representante).
  \item \textbf{Codificar} con \texttt{sentence-transformers}, modelo\\
  \texttt{paraphrase-multilingual-MiniLM-L12-v2} ($\approx$120\,MB, Hugging Face en la primera corrida). Vectores normalizados.
  \item \textbf{Bloquear por departamento.} Cusco no se compara con Cajamarca. Evita homónimos nacionales y reduce costo.
  \item \textbf{Clustering aglomerativo} (scikit-learn, métrica coseno, linkage average). Umbral de similitud por defecto \textbf{0{,}90}: se fusionan vectores cuya distancia coseno es $\le 1-0{,}90=0{,}10$.
  \item Asignar \texttt{caso\_id} entero a cada cluster; filas sin texto reciben id propio.
\end{enumerate}

Código de referencia: \texttt{02\_codigo/defensoria\_extract/identidad\_emb.py}.

\subsection{Por qué el umbral es 0{,}90}

Se calibró con el conflicto de \textbf{Cenepa} (Amazonas): la Defensoría describe oposición minera de las mismas comunidades frente a \textbf{Afrodita} y frente a \textbf{Dorato Perú}. Por debajo de 0{,}90 el modelo las fusionaba en un solo caso y borraba una distinción que importa para minería. El umbral alto \textbf{prefiere fragmentar} antes que mezclar empresas distintas.

Costo de ese criterio: conflictos muy reescritos (Espinar con varias fichas paralelas) pueden quedar en 2--3 \texttt{caso\_id} en vez de uno. Eso se audita a mano o se corrige después; no se ``arregla'' bajando el umbral a ciegas.

\subsection{Cómo se calibra y se audita}

\texttt{comparar\_identidad.py} no relee PDF. Contrasta el agrupamiento fuzzy contra embeddings a varios umbrales (p.\ ej.\ 0{,}82 / 0{,}86 / 0{,}90) y lista las \textbf{fusiones más grandes} para revisión humana. Salida: \texttt{comparacion\_identidad.txt}.

Regla práctica: un umbral bajo junta de más; uno alto parte de más. La decisión correcta es de contenido (¿misma empresa? ¿mismo pedido?), no de métrica sola.

\subsection{Alternativa fuzzy}

\begin{verbatim}
python -u 02_codigo/extraer_conflictos.py --identidad fuzzy
\end{verbatim}

Fuzzy por departamento, umbral 88 (\texttt{identidad.py}). Sigue disponible para comparar o reproducir corridas viejas. El default del proyecto es embeddings.

\subsection{Qué no hace la identidad automática}

\begin{itemize}
  \item No geocodifica ni nombra la unidad minera.
  \item No fusiona clusters que el humano ve como el mismo conflicto si los vectores quedaron lejos.
  \item No separa dos conflictos si el texto es casi idéntico y solo cambia un actor secundario (el umbral 0{,}90 mitiga el caso Afrodita/Dorato, no todos los casos límite).
\end{itemize}

%----------------------------------------------------------------------
\section{Cómo mejorar el resultado: trabajo manual y Gemini}

El panel ya es usable para conteos y series. Para un producto de \emph{investigación} más limpio (trayectorias largas sin huecos de id, nombres de mina, resúmenes legibles, precisión medida) hace falta una capa encima. Hay dos vías que se complementan: \textbf{manual} y \textbf{Gemini} (Vertex o API key).

\subsection{Qué sigue siendo manual (alto valor)}

\begin{enumerate}
  \item \textbf{Muestra dorada.} Generar con \texttt{generar\_muestra\_dorada.py}, anotar columnas \texttt{ok\_*} en el Excel, medir con \texttt{evaluar\_extraccion.py}. Sin eso no se sabe en qué maqueta falla el parseo.
  \item \textbf{Auditoría de identidad.} Revisar fusiones y fragmentaciones que lista \texttt{comparar\_identidad.py} (y casos largos como Espinar: varios \texttt{caso\_id} para el mismo territorio). Decidir fusiones/excepciones a mano.
  \item \textbf{Curación minera} (si el entregable es el tablero del proyecto). Sobre \texttt{menciona\_mineria=1}: confirmar mina, compañía, pueblos, ilegal/informal. El catálogo de Drive ayuda; el juicio es humano.
  \item \textbf{Basura residual.} Filas con caso muy corto, badges \texttt{NUEVO}/\texttt{REACTIVADO}, departamento vacío: borrar, fusionar o re-extraer ese PDF.
  \item \textbf{Inspección de maqueta 2024--26.} En reportes tardíos cae el número de fichas: ¿diseño Defensoría (solo activos con hechos) o corte del detector? Se confirma abriendo 2--3 PDF.
  \item \textbf{Junio 2005.} Conseguir el PDF faltante; el código no lo inventa.
\end{enumerate}

\subsection{Qué aporta Gemini (capa aditiva)}

\texttt{enriquecer\_llm.py} \textbf{no relee PDF} ni cambia el panel. Lee el panel, llama al modelo \textbf{una vez por \texttt{caso\_id}} (no por fila-mes) y escribe \texttt{resumenes\_llm.csv}.

Campos típicos del esquema JSON: resumen (2--3 oraciones), empresa, mineral, comunidades, demanda principal, tipo normalizado, confianza, texto ilegible.

\begin{itemize}
  \item Autenticación por defecto: \textbf{Vertex AI}, proyecto \texttt{researchassistant01}, ADC (\texttt{gcloud auth application-default login}).
  \item Alternativa: \textbf{API key} de Gemini Developer API en el \texttt{.env}:
\end{itemize}

\begin{verbatim}
GEMINI_AUTH=api_key
GEMINI_API_KEY=...
GEMINI_MODELO=gemini-2.5-flash
\end{verbatim}

(Ver \texttt{.env.ejemplo}.) Modelo por defecto: \texttt{gemini-2.5-flash}. Caché SQLite (\texttt{cache\_llm.sqlite}): repetir la corrida no vuelve a pagar casos ya hechos.

\begin{verbatim}
python -u 02_codigo/enriquecer_llm.py --limite 8
python -u 02_codigo/enriquecer_llm.py --solo-mineria
python -u 02_codigo/enriquecer_llm.py
\end{verbatim}

El CSV \texttt{casos\_unicos.csv} conserva el texto original. El Excel \texttt{casos\_unicos\_resumen.xlsx} puede usar el resumen de Gemini como columna legible.

\subsection{Dónde Gemini ayuda de verdad}

\begin{center}
\renewcommand{\arraystretch}{1.15}
\begin{tabularx}{\textwidth}{@{}l Y Y@{}}
\toprule
\textbf{Tarea} & \textbf{Gemini} & \textbf{Manual / código} \\
\midrule
Resumen legible del caso & Fuerte (una vez por id) & Costoso a escala \\
Empresa / mineral / comunidades & Útil como propuesta & Validar muestra; no inventar \\
Unir clusters fragmentados & Posible como sugerencia & Decisión final humana \\
Parsear latentes sin \texttt{Tipo:} & No sustituye un parser & Código + revisión PDF \\
Medir precisión de extracción & No & Muestra dorada \\
Decidir umbral de embeddings & No & \texttt{comparar\_identidad} + ojo \\
\bottomrule
\end{tabularx}
\end{center}

Reglas del prompt (ya en \texttt{llm/prompts.py}): usar \textbf{solo} el texto entregado; no inventar empresas ni lugares; vacío si no está en el insumo; marcar \texttt{texto\_ilegible} si la ficha es basura.

\subsection{Combinación recomendada}

\begin{enumerate}
  \item Dejar el panel como fuente de verdad del universo.
  \item Correr Gemini sobre minería (\texttt{--solo-mineria}) o sobre todos los ids, con Vertex o API key.
  \item Revisar a mano una muestra de resúmenes (confianza baja / texto ilegible / empresas dudosas).
  \item Usar la auditoría de identidad + curación de minas para el producto ``tablero del proyecto'', sin pisar el panel crudo.
\end{enumerate}

En corto: \textbf{embeddings} arman el universo de ids; \textbf{manual} calibra y cura; \textbf{Gemini} redacta y propone campos ---no reemplaza la extracción ni la muestra dorada.

%----------------------------------------------------------------------
\section{Qué pregunta responde cada producto}

\begin{center}
\renewcommand{\arraystretch}{1.15}
\begin{tabularx}{\textwidth}{@{}l Y@{}}
\toprule
\textbf{Producto} & \textbf{Pregunta} \\
\midrule
\texttt{panel\_caso\_mes.csv} & ¿Qué conflictos reportó la Defensoría en el mes $t$? \\
Filtro socioambiental / minería & ¿Cuántos de ese tipo hay en $t$ o en el departamento $d$? \\
\texttt{casos\_unicos} + trayectoria & ¿Cuánto dura este \texttt{caso\_id} en el registro? \\
\texttt{resumenes\_llm.csv} & ¿Cómo se cuenta el caso en dos oraciones? ¿Qué empresa cita el texto? \\
Tabla tipo Drive (catálogo) & ¿Cuáles son las historias de \emph{nuestras} minas? \\
\bottomrule
\end{tabularx}
\end{center}

El CSV de 83 de Drive responde la última pregunta. El pipeline local responde las primeras; la última se construye encima.

%----------------------------------------------------------------------
\section{Reproducir}

\begin{verbatim}
pip install -r requirements.txt
python -u 02_codigo/extraer_conflictos.py
python -u 02_codigo/comparar_identidad.py
python -u 02_codigo/generar_muestra_dorada.py
python -u 02_codigo/enriquecer_llm.py --solo-mineria
\end{verbatim}

Paquete por archivo: \texttt{layout.py}, \texttt{secciones.py}, \texttt{fichas.py}, \texttt{flags.py}, \texttt{identidad.py}, \texttt{identidad\_emb.py}. Capa LLM: \texttt{02\_codigo/llm/}.

%----------------------------------------------------------------------
\section{Cierre}

No se dejó el pipeline de Drive por capricho ni se inventó este por moda. El local prioriza \textbf{leer el PDF con geometría}, \textbf{no borrar el universo} y \textbf{agrupar por significado} (embeddings) con un umbral explícito y auditable. Lo que aún duele ---fragmentación de ids, latentes tabulares, nombres de mina, resúmenes--- se ataca con revisión humana y, donde escala, con Gemini vía Vertex o API key, siempre como capa encima del panel, no como sustituto del extractor.

\end{document}
